% Generalized Attention Flow (GAF): Feature Attribution via Maximum Flow
% Authors: Behrooz Azarkhalili, Maxwell Libbrecht (SFU, UC Berkeley)
% Consolidated technical content

\section{Abstract}
GAF (Generalized Attention Flow) extends Attention Flow by replacing attention weights with the generalized Information Tensor, integrating attention weights, their gradients, maximum flow problem, and barrier method. GAF satisfies key theoretical properties (Shapley axioms) and consistently outperforms 13 SOTA feature attribution methods on sequence classification.

\section{Problem with Existing Methods}
1. Gradient-based methods focus on gradients of attention weights, inheriting attention limitations
2. LRP approaches rely on specific rules limiting cross-architecture applicability
3. Most methods fail to capture pairwise feature interactions
4. Many violate axioms: symmetry, sensitivity, efficiency, linearity

\section{Key Issue: Non-Uniqueness in Attention Flow}
Previous work claimed Attention Flow produces Shapley values, but overlooked that maximum flow solutions are NOT unique. Non-unique solutions are frequent in practice and undermine theoretical properties.

\section{Methodology}

\subsection{Multi-Head Attention Representation}
Transformer attention weights:
\begin{equation}
\mathbf{A} = \text{Concat}(\hat{\mathbf{A}}_1, \ldots, \hat{\mathbf{A}}_l) \in \mathbb{R}^{l \times h \times t \times t}
\end{equation}
where $\hat{\mathbf{A}}_j \in \mathbb{R}^{h \times t \times t}$ is multi-head attention for layer $j$.

\subsection{Information Tensor}
Generalized information tensor $\bar{\mathbf{A}} \in \mathbb{R}^{l \times t \times t}$ aggregates attention weights. Three variants:

\textbf{1. Attention Flow (AF):}
$\bar{\mathbf{A}} := \mathbb{E}_h(\mathbf{A})$

\textbf{2. Gradient Flow (GF):}
$\bar{\mathbf{A}} := \mathbb{E}_h(\lfloor \nabla \mathbf{A} \rfloor_+)$
where $\nabla \mathbf{A} = \frac{\partial y_t}{\partial \mathbf{A}}$

\textbf{3. Attention-Gradient Flow (AGF):}
$\bar{\mathbf{A}} := \mathbb{E}_h(\lfloor \mathbf{A} \odot \nabla \mathbf{A} \rfloor_+)$

Here $\lfloor x \rfloor_+ = \max(x, 0)$ and $\odot$ is Hadamard product.

\subsection{Maximum Flow Problem}
Treat attention graph as flow network:
- Nodes: tokens across layers + super-source/super-target
- Edges: attention connections
- Capacities: values from information tensor

Minimum-cost circulation formulation:
\begin{equation}
\min_{\mathbf{B}^\top \mathbf{f} = 0, \mathbf{l} \leq \mathbf{f} \leq \mathbf{u}} \mathbf{c}^\top \mathbf{f}
\end{equation}

\subsection{Log Barrier Regularization (Key Innovation)}
To ensure UNIQUE solutions, add log barrier regularization:
\begin{equation}
\mathbf{f}^*_\mu = \arg\min_{\alpha(\mathbf{f})=0} \xi(\mathbf{f}) + \mu \psi(\mathbf{f})
\end{equation}
where $\psi(\mathbf{f}) = -\sum \log(-\beta(\mathbf{f}))$ is log barrier function.

\textbf{Theorem:} For strictly convex barrier $\psi$, convex $\xi$, and $\mu > 0$:
- There exists unique optimal point $\mathbf{f}^*_\mu$
- $\lim_{\mu \to 0} \mathbf{f}^*_\mu = \mathbf{f}^*$

\subsection{Feature Attribution Computation}
After solving regularized max-flow, feature attributions are flows from super-source to input tokens in first layer.

\section{Theoretical Properties}

GAF attributions are proven to be Shapley values satisfying:
1. \textbf{Efficiency:} Sum of attributions equals model output
2. \textbf{Symmetry:} Symmetric features get equal attribution
3. \textbf{Nullity:} Features with no contribution get zero attribution
4. \textbf{Linearity:} Attributions scale linearly with model outputs

These properties hold ONLY with barrier regularization ensuring uniqueness.

\section{Experiments}

\subsection{Datasets}
SST2, IMDB, Yelp, Amazon, AG News (sentiment and topic classification)

\subsection{Baselines (13 methods)}
- Gradient-based: Gradients, InputXGradient
- Game-theoretic: LIME, KernelSHAP
- Propagation: LRP variants
- Attention-based: Attention Rollout, Attention Flow, AttCAT
- Others: Integrated Gradients, DeepLIFT

\subsection{Evaluation Metrics}
\textbf{AOPC (Area Over Perturbation Curve):}
Measures prediction change when removing top-k important tokens

\textbf{LOdds (Log Odds):}
Measures log-odds change when masking important features

\subsection{Key Results}
AGF (Attention-Gradient Flow) consistently outperforms:
- All 13 baseline methods in most settings
- Original Attention Flow (no barrier) by significant margin
- Gradient-only methods (GF)

\textbf{AOPC Results (higher = better):}
| Dataset | Best Baseline | AGF |
|---------|--------------|-----|
| SST2 | ~0.35 | 0.42 |
| IMDB | ~0.30 | 0.38 |
| Yelp | ~0.32 | 0.40 |

\section{Key Contributions}
1. \textbf{Information Tensor:} Three ways to aggregate attention (AF, GF, AGF)
2. \textbf{Barrier Regularization:} Ensures unique max-flow solutions
3. \textbf{Shapley Guarantee:} Proves GAF attributions satisfy all Shapley axioms
4. \textbf{SOTA Performance:} AGF variant outperforms existing methods
5. \textbf{Open Source:} Flexible Python package for any HuggingFace encoder-only Transformer

\section{Why AGF Works Best}
- Attention weights capture feed-forward information flow
- Gradients capture backpropagation information flow
- Product captures BOTH simultaneously
- Max-flow naturally models feature interactions (unlike independent gradient computation)

\section{Computational Aspects}
- Build capacity network from information tensor
- Add super-source connected to first layer tokens
- Add super-target connected from last layer tokens
- Solve regularized max-flow problem
- Extract flows to input layer as attributions
